\documentclass[11pt, a4paper]{article}
\usepackage{geometry}
\usepackage[parfill]{parskip}
\usepackage{graphicx}
\usepackage{amsmath}
\usepackage{enumerate}

\usepackage{charter}

\title{Voorbeeldtentamen 3 \\ IB0402 Logica, verzamelingen en relaties}
\author{Harman Singh}
\date{Month Day, Year}

\begin{document}
\maketitle
% chktex-file 44

\pagebreak
\section*{Opgave 1}
\subsection*{oefening a}

\[
\begin{tabular}{|c|c|c|c|c|c|c|c|c|c|c|}
	\hline
	${p}$ & ${q}$ & $\neg{p}$ & ${p}\land{q}$ & ${p}\rightarrow{q}$ & ${p}\leftrightarrow{q}$ & ${p}\land{\neg{q}}$ & $\neg{({p}\land{\neg{q}})}$ & ${p}\lor{q}$  & $\neg{({p}\land\neg{q})} \rightarrow({p}\lor{q})$ \\
	\hline
	1 & 1 & 0 & 1 & 1 & 1 & 0 & 1 & 1 & 1 \\
	1 & 0 & 0 & 0 & 0 & 0 & 1 & 0 & 1 & 1 \\
	0 & 1 & 1 & 0 & 1 & 0 & 0 & 1 & 1 & 1 \\
	0 & 0 & 1 & 0 & 1 & 1 & 0 & 1 & 0 & 0 \\
	\hline
\end{tabular}
\]

We zoeken $\neg (p \land \neg{q}) \rightarrow (p \lor q)$

Hiervoor zoeken we best de volgende elementen:

\begin{itemize}
	\item $\neg{q}$
	\item $({p}\land{\neg{q}})$
	\item $\neg{({p}\land{\neg{q}})}$
\end{itemize}

\subsection*{oefening b}

Dit is geen tautologie (een uitspraak die altijd waar is) of een contradictie (een uitspraak die altijd onwaar is), maar een $\boxed{\text{contingentie}}$ omdat het waar (1) is in 3 van de 4 gevallen maar onwaar (0) in 1 van de 4 gevallen (namelijk waar $p$ en $q$ beiden 0 zijn)


\pagebreak
\section*{Opgave 2}

$\forall x P(x,y) \rightarrow \neg \exists y(Q(y) \lor R(y))$

$\equiv \forall x P(x,y) \rightarrow 		\neg \exists z(Q(z) \lor R(z)) $ \hfill (alfabetische variant)

$\equiv \forall x P(x,y) \rightarrow 		\forall z \neg(Q(z) \lor R(z)) $ \hfill (dualiteit)

$\equiv \neg \forall x P(x,y) \lor 		\forall z \neg(Q(z) \lor R(z)) $ \hfill ($\rightarrow$-eliminatie)

$\equiv \exists x \neg P (x,y) \lor 		\forall z \neg(Q(z) \lor R(z)) $ \hfill (dualiteit)

$\equiv \exists x \neg P (x,y) \lor 		\exists x \forall z \neg(Q(z) \lor R(z)) $ \hfill (loze existentiele kwantificatie)

$\equiv \exists x (\neg P (x, y)) \lor (\forall z  \neg(Q(z) \lor R(z)))$ \hfill (distributiviteit)

$\equiv \exists x (\forall z \neg P (x, y)) \lor (\forall z  \neg(Q(z) \lor R(z)))$ \hfill (loze universele kwantificatie)

$\equiv \exists x (\forall z \neg P (x, y)) \lor (\neg(Q(z) \lor R(z)))$ \hfill (loze universele kwantificatie)

$\equiv \boxed{\exists x \forall z (\neg P (x, y) \lor \neg(Q(z) \lor R(z)))}$



\pagebreak
\section*{Opgave 3}
\ldots



\pagebreak
\section*{Opgave 4}

${({A}^c \cup B)}^c \cap {(A \cap B)}^c$

$\equiv  ({(A^c)}^c \cap B^c) \cap ({A}^c \cup {B}^c)$ \hfill (De Morgan, beide kanten)

$\equiv  (A \cap B^c) \cap ({A}^c \cup {B}^c)$ \hfill (dubbel complement)

$\equiv  [(A \cap B^c) \cap {A}^c] \cup [(A \cap B^c) \cap {B}^c]$ \hfill (distributiewet)

$\equiv  [(A \cap A^c) \cap {B}^c] \cup [(A \cap B^c) \cap {B}^c]$ \hfill (associativiteit)

$\equiv  (\emptyset \cap {B}^c) \cup [(A \cap B^c) \cap {B}^c]$ \hfill (complementregels)

$\equiv  \emptyset \cup [(A \cap B^c) \cap {B}^c]$ \hfill (nulelement)

$\equiv  [(A \cap B^c) \cap {B}^c]$ \hfill (nulelement)

$\equiv  [A \cap (B^c \cap {B}^c)]$ \hfill (associativiteit)

$\equiv  [A \cap ({B}^c)]$ \hfill (idempotentie)

$\equiv \boxed{A \cap B^c}$



\section*{Opgave 5}

\subsection*{oefening a}
\ldots


\subsection*{oefening b}

$\boxed{\text{i en iii zijn equivalent}}$

$(\bar{x} \sqcap y) \sqcup \overline{(z \sqcup \overline{y})}$

$\equiv (\bar{x} \sqcap y) \sqcup (\bar{z} \sqcap \bar{\bar{y}})$ \hfill (De Morgan)

$\equiv (\bar{x} \sqcap y) \sqcup (\bar{z} \sqcap y)$ \hfill (dubbel complement)

$\equiv y \sqcap (\bar{x} \sqcup \bar{z})$ \hfill (distributiviteit, maar omgekeerd)

$\equiv y \sqcap \overline{(x \sqcap z)}$ \hfill (De Morgan)






\end{document}